%! Systems of Differential Equations — Unit 6, Lesson 6.4 — Homework (Answer Key)
\documentclass[10pt]{article}
\usepackage{linalg-article}
\usepackage{linalg-key}   % pulls in -boxes; do NOT also load -boxes
\newcommand{\TallMath}[1]{$\displaystyle #1\rule[-1.4em]{0pt}{3.2em}$}
\newcommand{\uu}{\mathbf{u}}
\newcommand{\xx}{\mathbf{x}}
\newcommand{\zero}{\mathbf{0}}
\newcommand{\T}{\mathsf{T}}

\begin{document}

\pageheader{Unit 6, Lesson 6.4}{Homework --- Answer Key}
\namedateperiod

\vspace{0.12in}

\begin{notesbox}{Practice}
To solve $\frac{d\uu}{dt}=A\uu$: split $\uu(0)=c_1\xx_1+c_2\xx_2$ into eigenvectors, then
$\uu(t)=c_1e^{\lambda_1 t}\xx_1+c_2e^{\lambda_2 t}\xx_2$. The single equation $\frac{du}{dt}=\lambda u$ has
$u(t)=u(0)e^{\lambda t}$ --- grows if $\lambda>0$, decays if $\lambda<0$.

\begin{enumerate}[leftmargin=1.9em, itemsep=12pt]
  \item \textbf{Single equations.} Solve each, then say grow or decay.
    \begin{enumerate}[label=(\alph*), leftmargin=1.8em, itemsep=4pt]
      \item $\frac{du}{dt}=4u$, \ $u(0)=3$ \hfill \ans{$u(t)=3e^{4t}$ --- grows.}
      \item $\frac{dw}{dt}=-3w$, \ $w(0)=10$ \hfill \ans{$w(t)=10e^{-3t}$ --- decays.}
    \end{enumerate}
  \item \textbf{General solution.} For $A=\begin{bmatrix} 1 & 2 \\ 2 & 1 \end{bmatrix}$
        ($\lambda=3,\xx_1=\begin{bsmallmatrix}1\\1\end{bsmallmatrix}$; $\lambda=-1,\xx_2=\begin{bsmallmatrix}1\\-1\end{bsmallmatrix}$),
        write the general solution of $\frac{d\uu}{dt}=A\uu$.
        \ansline{$\uu(t)=c_1e^{3t}\begin{bsmallmatrix}1\\1\end{bsmallmatrix}+c_2e^{-t}\begin{bsmallmatrix}1\\-1\end{bsmallmatrix}$.}
  \item \textbf{Fit the start.} Same $A$, starting at $\uu(0)=\begin{bsmallmatrix}4\\2\end{bsmallmatrix}$. Find
        $c_1,c_2$ and write the full $\uu(t)$. Which eigen-direction wins as $t\to\infty$?
        \ansline{$c_1+c_2=4,\ c_1-c_2=2\Rightarrow c_1=3,c_2=1$, so $\uu(t)=3e^{3t}\begin{bsmallmatrix}1\\1\end{bsmallmatrix}+1e^{-t}\begin{bsmallmatrix}1\\-1\end{bsmallmatrix}$. The $e^{-t}$ piece fades, so $\xx_1=\begin{bsmallmatrix}1\\1\end{bsmallmatrix}$ (the $\lambda=3$ direction) wins.}
  \item \textbf{Stability screen.} Each system $\frac{d\uu}{dt}=A\uu$ has the eigenvalues shown. Decide whether
        every solution \emph{decays to $\zero$}, \emph{grows}, or is \emph{mixed} (one of each).
    \begin{enumerate}[label=(\alph*), leftmargin=1.8em, itemsep=4pt]
      \item $\lambda=-2,\ -5$ \hfill \ans{decays to $\zero$ (stable)}
      \item $\lambda=3,\ 1$ \hfill \ans{grows}
      \item $\lambda=2,\ -4$ \hfill \ans{mixed --- grows overall}
    \end{enumerate}
    \ansline{(a) Both negative $\Rightarrow$ both modes decay $\Rightarrow\uu(t)\to\zero$. (b) Both positive $\Rightarrow$ both modes grow. (c) One of each: the $\lambda=2$ mode grows while $\lambda=-4$ decays, so the solution ultimately \emph{grows} along the $\lambda=2$ eigen-direction (a saddle).}
  \item \textbf{Explain.} (a) Why does a solution that starts on an eigenvector stay on that eigen-line?
        (b) Why do the signs of the eigenvalues decide the long-term behavior?
        \ansline{(a) On an eigenvector, $A\xx=\lambda\xx$, so $\frac{d\uu}{dt}=A\uu$ becomes the single equation $\frac{d\uu}{dt}=\lambda\uu$ --- pure scaling by $e^{\lambda t}$, which never leaves the line through $\xx$. (b) Each piece carries a factor $e^{\lambda t}$: if $\lambda>0$ it grows without bound, if $\lambda<0$ it decays to $0$. So the signs alone say which modes survive and which vanish.}
\end{enumerate}
\end{notesbox}

\begin{extensionbox}
\textbf{Running time forward, matrix-style.} Diagonalizing gives $\uu(t)=Xe^{\Lambda t}X^{-1}\uu(0)$, where
$e^{\Lambda t}=\begin{bsmallmatrix}e^{\lambda_1 t}&0\\0&e^{\lambda_2 t}\end{bsmallmatrix}$ --- the same shape as
$A^k=X\Lambda^k X^{-1}$, but with $\lambda^k$ replaced by $e^{\lambda t}$. For the Problem~2 matrix
($\lambda=3,-1$), write $e^{\Lambda t}$, and in one sentence explain why this reproduces the combination
$\uu(t)=c_1e^{3t}\xx_1+c_2e^{-t}\xx_2$.
\ansline{$e^{\Lambda t}=\begin{bsmallmatrix}e^{3t}&0\\0&e^{-t}\end{bsmallmatrix}$. Writing $X^{-1}\uu(0)=\begin{bsmallmatrix}c_1\\c_2\end{bsmallmatrix}$ (the eigenvector coordinates of the start), $e^{\Lambda t}$ scales each coordinate by its own $e^{\lambda t}$, and multiplying by $X$ rebuilds them as $c_1e^{3t}\xx_1+c_2e^{-t}\xx_2$.}
\end{extensionbox}

\begin{spiralbox}
\textbf{Unit 6 in one line.} You met eigenvalues ($A\xx=\lambda\xx$), used them to diagonalize
($A=X\Lambda X^{-1}$), saw the symmetric case ($S=Q\Lambda Q^{\T}$), and now let them run a system forward in
time ($e^{\lambda t}$). The same eigenvalues answer every question. \textbf{Next: the Unit~6 review and test.}
\end{spiralbox}

\begin{teachernote}
Item~1 is the scalar building block (grow vs.\ decay from the sign). Item~2 assembles the general solution on
the spine $\begin{bsmallmatrix}1&2\\2&1\end{bsmallmatrix}$; Item~3 fits $\uu(0)=\begin{bsmallmatrix}4\\2\end{bsmallmatrix}$
($c_1=3,c_2=1$) and reads the winning mode. Item~4 is a stability \emph{screen} --- (a) all-negative stable,
(b) all-positive growth, (c) the mixed ``saddle'' where the positive eigenvalue dominates. Item~5 is the two
target justifications ($A\xx=\lambda\xx$ collapses to the scalar equation; $e^{\lambda t}$ signs decide
survival). The extension ties $\uu(t)=Xe^{\Lambda t}X^{-1}\uu(0)$ back to $A^k=X\Lambda^k X^{-1}$ (replace
$\lambda^k$ by $e^{\lambda t}$). Common slips: dropping coefficients or the $e^{\lambda t}$ factor; believing a
positive and negative eigenvalue cancel; and reading ``mixed'' as ``stays put'' rather than ``grows.''
\end{teachernote}

\end{document}
